\subsection{Lines, Segments, and Rays} 
If we have two points, we can draw a line, ray, or segment.
\begin{itemize}
\item A line \makeblank[3.5in] in both directions.
\item A ray \makeblank[3.5in] in one direction, and \makeblank[3.5in] in the other.
\item A segment \makeblank[3.25in] in both directions.
\end{itemize}
We name these figures using the two points and a marking over them.
\begin{center}\begin{tabular}{ccc}
Symbol&Read&Image\\\hline
\fitrow{0.5in}{7/3}{$\overleftrightarrow{AB}$}&
\fitrow{2.5in}{7/3}{\makeblank}&
\begin{tikzpicture}
\useasboundingbox (-1,-2/3) rectangle (4,5/3);
\fill (0,0) circle (.1) node[anchor=north]{$A$};
\fill (3,1) circle (.1) node[anchor=north]{$B$};
\draw[<->,geom] (-1,-1/3) -- (4,4/3);
\end{tikzpicture}\\
\fitrow{0.5in}{7/3}{$\overleftrightarrow{BA}$}&
\fitrow{2.5in}{7/3}{\makeblank}&
\begin{tikzpicture}
\useasboundingbox (-1,-2/3) rectangle (4,5/3);
\fill (0,0) circle (.1) node[anchor=north]{$A$};
\fill (3,1) circle (.1) node[anchor=north]{$B$};
\draw[<->,geom] (-1,-1/3) -- (4,4/3);
\end{tikzpicture}\\
\fitrow{0.5in}{7/3}{$\overrightarrow{AB}$}&
\fitrow{2.5in}{7/3}{\makeblank}&
\begin{tikzpicture}
\useasboundingbox (-1,-2/3) rectangle (4,5/3);
\fill (0,0) circle (.1) node[anchor=north]{$A$};
\fill (3,1) circle (.1) node[anchor=north]{$B$};
\draw[->,geom] (0,0) -- (4,4/3);
\end{tikzpicture}\\
\fitrow{0.5in}{7/3}{$\overrightarrow{BA}$}&
\fitrow{2.5in}{7/3}{\makeblank}&
\begin{tikzpicture}
\useasboundingbox (-1,-2/3) rectangle (4,5/3);
\fill (0,0) circle (.1) node[anchor=north]{$A$};
\fill (3,1) circle (.1) node[anchor=north]{$B$};
\draw[<-,geom] (-1,-1/3) -- (3,1);
\end{tikzpicture}\\
\fitrow{0.5in}{7/3}{$\overline{AB}$}&
\fitrow{2.5in}{7/3}{\makeblank}&
\begin{tikzpicture}
\useasboundingbox (-1,-2/3) rectangle (4,5/3);
\fill (0,0) circle (.1) node[anchor=north]{$A$};
\fill (3,1) circle (.1) node[anchor=north]{$B$};
\draw[geom] (0,0) -- (3,1);
\end{tikzpicture}\\
\fitrow{0.5in}{7/3}{$\overline{BA}$}&
\fitrow{2.5in}{7/3}{\makeblank}&
\begin{tikzpicture}
\useasboundingbox (-1,-2/3) rectangle (4,5/3);
\fill (0,0) circle (.1) node[anchor=north]{$A$};
\fill (3,1) circle (.1) node[anchor=north]{$B$};
\draw[geom] (0,0) -- (3,1);
\end{tikzpicture}
\end{tabular}\end{center}
We can also denote a line with a \makeblank. In the figure below, \makeblanksmall is also called \makeblanksmall (read ``\makeblank[1.5in]'').
\begin{icpic}
\fill (0,0) circle (.1) node[anchor=north]{$P$};
\fill (2.5,.5) circle (.1) node[anchor=north]{$Q$};
\draw[<->,geom] (-1.5,-0.3) -- (6.5,1.3) node[anchor=south east]{$\ell$};
\end{icpic}